% %%%%%%%%%%%%%%%%%%%%%%%%%%%%%%%%%%%%%%%%%%%%%%%%%%%%%%%%%%%%
% CHAPTER 1 - INTRODUCTION
% %%%%%%%%%%%%%%%%%%%%%%%%%%%%%%%%%%%%%%%%%%%%%%%%%%%%%%%%%%%%
%
% DRAFT. Before submission:
%  - Read each cited paper and add page numbers to the citations
%    (SRH Guidelines 4.2.1), e.g. \parencite[p.~538]{gregory2021}.
%  - Rewrite the text in your own words and check every claim against the source.
%  - Confirm the scope decisions in Section 1.4 with your supervisor.

\chapter{Introduction}
\label{ch:intro}

\section{Background to the Research and Research Justification}
\label{sec:background}

Artificial intelligence (AI) has moved from a specialist technology to a general-purpose resource that can augment or replace human tasks across industries, including retail, logistics and supply chains \parencite{dwivedi2021}. E-commerce is one of the fields in which this shift is most visible, and research on AI applications in online commerce has grown into a substantial body of work in its own right \parencite{bawack2022}. Platform-based businesses are particularly exposed to this development. Many of the most profitable firms own platforms that exhibit network effects, and advances in AI have added a further mechanism, data network effects, through which a platform becomes more valuable to each user the more it learns from the data it collects \parencite{gregory2021}. Machine learning can therefore strengthen the lock-in effects of a business model beyond the traditional direct and indirect network effects \parencite{haftor2021}.

This logic has an important implication for competition. If more data leads to more learning, and more learning attracts more users who generate more data, the resulting cycle favours firms that are already large \parencite{haftor2021}. AI would then reinforce existing positions rather than create new ones. The relationship is, however, not automatic. \textcite{gregory2021} argue that the link between a platform's AI capability and the value perceived by its users depends on platform legitimation, data stewardship and user-centric design. Whether AI produces an advantage therefore depends on how the platform organises and governs its use of AI, and not only on the volume of data it controls. This distinction between owning AI and organising around it is the starting point of this thesis.

Investment in AI does not translate automatically into superior performance. A systematic review of the AI business value literature concludes that organisations still struggle to adopt and leverage AI in their operations, and that the mechanisms through which AI creates value remain insufficiently understood \parencite{enholm2022}. The same review separates first-order effects of AI, such as process efficiency, from second-order effects at the level of the organisation, a distinction that is central to the question of how AI might lead to a lasting competitive advantage. Comparable difficulties have been reported when manufacturing firms attempt to scale AI capabilities across their business models and operations \parencite{sjodin2021}. AI can even have negative consequences: opaque, poorly governed AI-integrated analytics have been associated with suboptimal decisions, operational inefficiency and, ultimately, competitive disadvantage \parencite{rana2022}. In that study, the opacity resulted from a lack of governance, poor data quality and insufficient training of key employees, which are organisational rather than technical shortcomings.

The positive evidence should also be read with care. An analysis of 500 AI deployment cases suggests that organisations achieve performance gains from AI capabilities only when they use them to reconfigure their processes \parencite{wambataguimdje2020}. The finding is relevant because it points to reconfiguration, a core element of dynamic capabilities, as the condition under which AI creates value. The cases were, however, collected from the websites of AI solution providers, a source that is likely to over-represent successful deployments. The study therefore indicates what AI can achieve under favourable conditions rather than how often such outcomes occur.

The resource-based view (RBV) offers one explanation for these differences. From this perspective, value does not stem from AI technology as such but from an AI capability, that is, a firm-specific combination of AI-related tangible, human and intangible resources \parencite{mikalef2021ai}. Empirical work supports a link between such a capability and firm outcomes, including organisational creativity and performance in the same study and, in a survey of 394 e-commerce firms, an indirect effect on performance through creativity, AI management and AI-driven decision making \parencite{chen2022}. In the latter study, the strength of these relationships also depended on contextual conditions such as innovation culture and environmental dynamism. AI may, moreover, change the sources of competitive advantage themselves. Using chess tournaments as a controlled setting, \textcite{krakowski2023} show that AI can make established capabilities obsolete while creating new sources of heterogeneity at the intersection of human and machine decision making. The advantage then lies neither in the human nor in the technology alone, but in how the two are combined.

Dynamic capabilities theory extends this reasoning to fast-changing environments. It shifts the attention from the resources a firm owns to its ability to sense changes, seize opportunities and transform its resource base. Case evidence from three large Norwegian organisations indicates that AI can act as a micro-foundation of such capabilities, helping firms to manage their operations in dynamic and uncertain environments \parencite{mikalef2021b2b}. Conceptual work similarly argues that AI leads to competitive advantage only when it is \enquote{situated} in the firm through grounding, bounding and recasting activities, which make AI-driven capabilities firm-specific, cost-effective and suited to the opportunities in the strategic environment \parencite{kemp2024}. Taken together, these contributions suggest that dynamic capabilities may be the mechanism that connects an AI capability with an advantage that lasts.

This mechanism deserves particular attention in e-commerce and platform-based business. Competition in digital markets follows a logic that departs in important respects from the predictions of mainstream competitive theory, because the value of a platform depends both on its size and on its distinctive technological and market identity \parencite{cennamo2021}. Findings obtained in manufacturing or in traditional business-to-business settings can therefore not simply be transferred to these firms. A clearer understanding of how AI capabilities contribute to a lasting competitive advantage in this context is relevant for two groups. Managers need to decide whether to invest primarily in AI technology or in the organisational capabilities required to use it, and researchers need to explain why similar AI investments lead to different competitive outcomes.

\section{Research Problem}
\label{sec:problem}

As AI tools, cloud services and pre-trained models become widely accessible, the technology alone offers limited protection against imitation. Rival firms can acquire similar software and hire similar specialists, so an advantage is more likely to depend on how AI is embedded in the organisation than on access to the technology itself \parencite{kemp2024, krakowski2023}. The central question is therefore not whether firms use AI, but through which organisational mechanisms an AI capability is converted into an advantage that persists over time.

The existing empirical evidence answers this question only in part. Several studies report a positive relationship between AI capability, or related constructs such as AI assimilation, and firm performance. They do so, however, through a wide range of intervening variables, including organisational creativity \parencite{mikalef2021ai, chen2022}, business-to-business marketing capabilities \parencite{mikalef2023}, organisational and customer agility \parencite{fossowamba2022} and business model innovation \parencite{shao2026}. Some of these mediators, such as agility or business model innovation, are conceptually related to sensing, seizing and transforming, but they are not identical with them. Only a few studies model the three dimensions of dynamic capabilities as distinct links between AI and competitive advantage. One recent study of 243 Spanish organisations and ten case studies relates AI integration to these dimensions and reports an overall indirect effect of AI-enabled dynamic capabilities on competitive advantage \parencite{selgascors2026}. Separate indirect effects for each dimension are not reported, so it remains open which of the three dimensions carries the effect.

A second difficulty concerns the constructs themselves. The literature does not use a single definition of AI capability. \textcite{mikalef2021ai} define it as a combination of AI-specific resources and treat it as a resource in the sense of the RBV. \textcite{chen2022} model it as a higher-order construct formed by basic resources, proclivity and skills, and \textcite{shao2026} distinguish tangible, intangible and skill resources. Other studies examine AI assimilation \parencite{fossowamba2022} or AI integration \parencite{selgascors2026}, while adoption studies focus on the earlier decision to introduce AI at all \parencite{chatterjee2021}. The conceptual status of AI capability is also contested: Shao and colleagues interpret it as a dynamic capability in its own right, whereas the RBV-based studies treat it as a resource that dynamic capabilities act upon. If AI capability is itself a dynamic capability, the proposed mediation through sensing, seizing and transforming partly overlaps with the independent variable. The thesis therefore needs to establish which construct each study actually measures before its findings can be compared.

The outcome variables are equally diverse. The studies cited above measure organisational creativity, firm performance, business-to-business marketing outcomes, competitive advantage and, in one case, competitive disadvantage \parencite{mikalef2021ai, mikalef2023, shao2026, rana2022}. Performance and competitive advantage are related but distinct concepts: a firm can perform well in a given year without holding an advantage that rivals cannot imitate. Treating performance as a proxy for sustainable competitive advantage would therefore overstate what the evidence shows.

The evidence is also concentrated in contexts that differ from e-commerce and platform-based business. Much of it comes from manufacturing firms \parencite{sjodin2021, chatterjee2021}, business-to-business marketing \parencite{mikalef2021b2b, mikalef2023}, small and medium-sized enterprises in China \parencite{shao2026}, service firms in India \parencite{rana2022} or cross-industry samples \parencite{selgascors2026}. Where e-commerce firms are studied directly, the outcome is general firm performance rather than a lasting advantage \parencite{chen2022}, while the work that explicitly addresses platforms is largely conceptual \parencite{gregory2021, cennamo2021}. The research designs vary as well, ranging from surveys analysed with structural equation modelling to qualitative case studies and, in one case, a controlled setting outside business \parencite{krakowski2023}. The survey-based studies typically rely on cross-sectional data and on outcome measures such as performance, which cannot by themselves establish whether an advantage persists over time or in which direction the relationships run.

It therefore remains unclear to what extent the published evidence supports the proposition that AI capabilities lead to a sustainable competitive advantage in e-commerce and platform-based businesses through sensing, seizing and transforming capabilities. A further study that simply adds another cross-sectional survey would not resolve this uncertainty, because the main obstacle is the fragmentation of the existing evidence. This thesis addresses the problem by systematically assessing the published empirical evidence for each of the proposed relationships, with explicit attention to how constructs were measured, in which context and with which research design.

\section{Research Objectives and Research Questions}
\label{sec:objectives}

The overall aim of this thesis is to evaluate the extent to which published empirical studies support the proposition that AI capabilities contribute to a sustainable competitive advantage in e-commerce and platform-based businesses through dynamic capabilities. The evaluation takes account of differences in construct measurement, context and research design between the studies. To achieve this aim, the thesis pursues six objectives. The first two concern the direct relationships on which the model rests; the remaining four concern the proposed mediating mechanisms, which form the main contribution of the thesis:

\begin{enumerate}
  \item To evaluate the evidence on the relationship between AI capabilities and sustainable competitive advantage.
  \item To evaluate the evidence on the relationship between AI capabilities and dynamic capabilities.
  \item To evaluate the evidence on the mediating role of sensing capability.
  \item To evaluate the evidence on the mediating role of seizing capability.
  \item To evaluate the evidence on the mediating role of transforming capability.
  \item To evaluate the evidence on the mediating role of dynamic capabilities as a higher-order construct.
\end{enumerate}

Objectives 3 to 5 and Objective 6 are not simply repetitions of the same question. The separate assessment of sensing, seizing and transforming shows where in the organisation AI creates value, which matters for managers who must decide where to direct their investment. The higher-order view in Objective 6 asks whether the three dimensions work as a combined capability whose effect is greater than that of each dimension alone. Because both perspectives may draw on the same studies, the evidence is recorded once per study and then assigned to the relevant objectives, so that no study is counted twice in support of the same claim. These objectives lead to four research questions:

\begin{description}
  \item[RQ1:] To what extent does published empirical evidence support a relationship between AI capabilities and sustainable competitive advantage?
  \item[RQ2:] To what extent does published empirical evidence support a relationship between AI capabilities and dynamic capabilities?
  \item[RQ3:] To what extent does published empirical evidence indicate that sensing, seizing and transforming capabilities each mediate the relationship between AI capabilities and sustainable competitive advantage?
  \item[RQ4:] To what extent does published empirical evidence indicate that dynamic capabilities, as a higher-order construct, mediate this relationship?
\end{description}

Chapter~\ref{ch:litreview} translates these questions into six hypotheses (H1--H6). Because the thesis is based on a structured review of existing studies, the hypotheses are assessed against the reported evidence rather than tested statistically with new data. A hypothesis can therefore be judged as supported, partly supported, inconclusive, contradicted or not yet examined by the available studies.

\section{Scope and Limitations}
\label{sec:scope}

The unit of analysis is the firm. The thesis focuses on e-commerce retailers and on firms that own and operate digital platforms. Complementors, such as independent sellers or app developers who offer their products through a platform, face different conditions regarding data access and governance. Studies of complementors are therefore treated as closely related evidence: they are counted as part of the wider e-commerce and platform context, but they are labelled as such and not taken as evidence about platform owners. Studies from other industries are included where they examine the proposed relationships, but their transferability to e-commerce and platform-based business is discussed explicitly rather than assumed.

In this thesis, a sustainable competitive advantage denotes an advantage over competitors that persists over time and resists imitation. The term does not refer to environmental sustainability. The distinction is necessary because parts of the related literature use the word in its ecological sense; \textcite{haftor2021}, for example, examine how machine learning can increase profitability while reducing negative ecological impacts. Studies that use \enquote{sustainable} only in this environmental sense are not treated as evidence on the persistence of competitive advantage.

Several questions lie outside the scope of the thesis. It does not evaluate individual AI technologies or algorithms, does not examine the consumer perspective on AI in e-commerce and does not measure the capabilities of particular firms. It also does not make causal claims beyond those that the underlying studies can support. Where the evidence is observational, relationships are described as associations rather than effects.

Further limitations follow from the research design. The thesis does not collect primary data and can only be as strong as the published studies on which it is based. Many of these studies are cross-sectional, rely on managers' perceptions and examine organisations outside e-commerce, so their findings are transferable to the focal setting only to a limited extent. Finally, the review is restricted to studies that can be identified through the documented search procedure; the exact databases, time window and eligibility criteria are specified in Chapter~\ref{ch:method}.

\section{Research Approach}
\label{sec:approach}

The thesis follows a deductive approach. Hypotheses are derived from the resource-based view and dynamic capabilities theory in Chapter~\ref{ch:litreview} and then compared with the findings of existing empirical studies. A structured literature review with a narrative synthesis is used for this purpose \parencite{booth2025}. This design suits the research questions for three reasons. First, the main uncertainty identified in Section~\ref{sec:problem} lies in the fragmentation of existing evidence, which a further single study could not resolve. Second, the relevant studies measure AI capability, dynamic capabilities and competitive advantage in different ways, which makes a statistical pooling of their results in a meta-analysis inappropriate. Third, original survey data from e-commerce and platform firms would be difficult to obtain within the scope of a master's thesis and would again be limited to a single, cross-sectional context.

The review proceeds in four steps. First, relevant studies are identified through a documented and reproducible database search. Second, explicit inclusion, exclusion and quality criteria are applied. The quality appraisal considers, among other aspects, whether a study separates its measures in time, how it addresses common method bias and whether indirect effects were actually estimated rather than inferred. Third, the evidence from each included study is recorded in an extraction matrix that distinguishes direct associations from reported indirect effects and documents the exact constructs measured. Fourth, the evidence for each hypothesis is classified as directly reported, indirectly relevant, contradictory or not available, taking the methodological quality of each study into account.

This approach has a clear limitation: it cannot produce new estimates of the proposed relationships. Its contribution lies instead in establishing, transparently and for each relationship separately, how much the current evidence can support, and in identifying where further empirical research is most needed.

\section{Structure of the Thesis}
\label{sec:structure}

The thesis consists of five chapters. Following this introduction, Chapter~\ref{ch:litreview} reviews the theoretical foundations of the resource-based view and dynamic capabilities theory, defines the key constructs and examines the empirical literature on AI capability in e-commerce and platform-based business. It closes by identifying the research gap and developing the conceptual model and hypotheses. Chapter~\ref{ch:method} then describes the research design, including the search strategy, the eligibility criteria, the quality appraisal and the synthesis rules.

Chapter~\ref{ch:analysis} presents the results of the review. It reports the characteristics of the included studies, assesses the evidence for each hypothesis and discusses the findings in relation to the literature reviewed in Chapter~\ref{ch:litreview}. Chapter~\ref{ch:conclusion} summarises the findings, answers the research questions and discusses the theoretical and managerial implications, the limitations of the thesis and directions for future research.
